To assess whether depth signal extends beyond handcrafted features, we fine-tune a pretrained language model end-to-end as a depth classifier using QLoRA \cite{qlora}.

\paragraph{Setup.}
We fine-tune Qwen3-0.6B-Base with LoRA rank 16, learning rate $2 \times 10^{-4}$, and 3 epochs, framing depth estimation as 5-class sequence classification (gen-0 through gen-4).\footnote{Evaluated on a 4-generator subset (Pythia-1B, OLMo-2-1B, Qwen3-0.6B, Qwen3-1.7B) using 5-fold stratified CV; XGBoost 3-class F1\,=\,0.805 under the same protocol. We expect 8-generator training to further improve given additional diversity; the comparison direction is preserved.}
For 3-class evaluation, we collapse predictions into gen-0\,/\,gen-1\,/\,gen-2+ post hoc.
All results use 5-fold stratified cross-validation on the 4-generator subset.

\begin{table}[ht]
\centering
\small
\caption{QLoRA end-to-end depth classification vs.\ XGBoost on handcrafted features, under 5-fold stratified CV (4-generator subset).}
\label{tab:qlora}
\resizebox{\linewidth}{!}{%
\begin{tabular}{lccc}
\toprule
\textbf{Method} & \textbf{5-class F1} & \textbf{3-class F1} & \textbf{Per-class F1 (5-class)} \\
\midrule
XGBoost (7 features) & 0.497 $\pm$ 0.008 & 0.805 $\pm$ 0.008 & --- \\
QLoRA (Qwen3-0.6B)   & 0.583 $\pm$ 0.010 & 0.905 $\pm$ 0.003 & 0.999\,/\,0.796\,/\,0.413\,/\,0.259\,/\,0.447 \\
\bottomrule
\end{tabular}}
\end{table}

\paragraph{Results.}
QLoRA achieves 3-class F1\,=\,0.905 ($\pm$0.003) under 5-fold stratified CV on a 4-generator subset, a +10.0\,pp improvement over the same-subset XGBoost baseline (0.805).
On 5-class evaluation, QLoRA reaches F1\,=\,0.583 ($\pm$0.010) vs.\ 0.497 for XGBoost (+8.6\,pp), with the largest gains at gen-1 (0.796 vs.\ 0.650) and gen-2 (0.413 vs.\ 0.267).
The gen-2/gen-3 boundary remains the primary failure mode (gen-3 F1\,=\,0.259), consistent with the feature-level analysis in \S\ref{sec:discussion}.

\paragraph{Hyperparameter robustness.}
A second run with LoRA rank 32 and seed 123 yields 5-class F1\,=\,0.587 ($\pm$0.005) and 3-class F1\,=\,0.905 ($\pm$0.004), which confirms stability across hyperparameter and seed variation ($\Delta < 1$\,pp on both metrics).

\paragraph{Implications.}
The substantial 3-class improvement demonstrates that raw token embeddings contain depth signal beyond what summary statistics capture; the seven handcrafted features are a strong but lossy compression of the full distributional information.
End-to-end fine-tuning exploits token-level patterns (e.g., positional regularities, n-gram repetitions) that aggregate features cannot represent.
This confirms that the depth signal documented in our main analysis is real and exploitable beyond engineered features, while the handcrafted approach retains practical advantages: interpretability, computational efficiency (no GPU inference), and the ability to reveal the differential depth response pattern that motivates the task formulation.
The QLoRA ceiling thus represents an information-asymmetric upper bound rather than a competing method: it consumes the full token sequence (0.6B parameters, GPU inference) whereas the 7-feature ensemble distills the signal into interpretable scalars that expose \emph{why} depth is detectable, a decomposition that end-to-end models cannot provide.
